% Irodov, Arihant PDF pp. 15--16; original printed pages 6--7.
% Problems 31--40. Derivations follow source data, not blindly the printed key.

\exproblem{31}{6}
一小球从静止开始竖直落下，落下高度 $h$ 后与倾角为 $\alpha$ 的光滑斜面
发生完全弹性碰撞。求再次碰到该斜面的位置与首次碰撞点沿斜面的距离。
\exsolution
首次撞击前球速 $v=\sqrt{2gh}$，方向竖直向下。
以沿斜面向下为切向、离开斜面为法向，碰撞前切向速度为 $v\sin\alpha$，
法向速度为 $-v\cos\alpha$。光滑弹性碰撞保持切向速度、反转法向速度，
所以碰后切向和法向初速分别为 $v\sin\alpha$、$v\cos\alpha$。
重力在两个方向上的加速度分别为 $g\sin\alpha$ 和 $-g\cos\alpha$。
再次撞面要求法向位移为零：
\[
 0=v\cos\alpha\,T-\frac12g\cos\alpha\,T^2
 \quad\Longrightarrow\quad T=\frac{2v}{g}.
\]
此时沿斜面位移
\[
 \ell=v\sin\alpha\,T+\frac12g\sin\alpha\,T^2
 =\frac{4v^2\sin\alpha}{g}.
\]
故
\[
 \boxed{\ell=8h\sin\alpha.}
\]
\excheck 取 $\alpha=0$ 得 $\ell=0$，符合水平镜面反弹后落回原地的极限。
此题的“落下高度”按竖直下落的高度理解。

\exproblem{32}{6}
火炮与目标位于同一高度，相距 $L=5.10\,\mathrm{km}$。
炮弹初速 $v_0=240\,\mathrm{m/s}$，不计阻力。问发射后多久击中目标？
\exsolution
同高射程满足
\[
 L=\frac{v_0^2\sin2\theta}{g},\quad
 q=\frac{gL}{v_0^2}\simeq0.8686<1.
\]
因此有两个可行仰角
\[
 \theta_1=\tfrac12\arcsin q\simeq30.15^\circ,\quad
 \theta_2=\tfrac12(\pi-\arcsin q)\simeq59.85^\circ .
\]
弹道飞行时间为 $T=L/(v_0\cos\theta)$，代入分别得到
\[
 \boxed{T_1\simeq24.6\,\mathrm s=0.410\,\mathrm{min},\qquad
 T_2\simeq42.3\,\mathrm s=0.705\,\mathrm{min}.}
\]
\excheck 必须验证 $L\le v_0^2/g$ 才有可行同高弹道。
书末给出的约 $0.41$ 与 $0.71\,\mathrm{min}$ 对应低、高两条弹道。

\exproblem{33}{6}
同一火炮在同一方位角内，以共同初速 $v_0=250\,\mathrm{m/s}$
先后发射两枚炮弹。第一枚仰角 $\theta_1=60^\circ$，
第二枚仰角 $\theta_2=45^\circ$。忽略阻力，求使两弹空中相撞的发射间隔。
\exsolution
设碰撞发生时两枚炮弹已飞行 $t_1,t_2$，其中 $t_1>t_2>0$，
发射间隔 $\Delta=t_1-t_2$。水平位置重合意味着
$t_1\cos\theta_1=t_2\cos\theta_2$；竖直位置重合意味着
\[
 v_0(t_1\sin\theta_1-t_2\sin\theta_2)
 =\frac g2(t_1^2-t_2^2).
\]
由前式置 $t_1=k\cos\theta_2,\ t_2=k\cos\theta_1$，
可将竖直位置方程改写成
\[
 v_0k\sin(\theta_1-\theta_2)
 =\frac g2 k^2(\cos^2\theta_2-\cos^2\theta_1).
\]
由于 $\cos^2\theta_2-\cos^2\theta_1
=(\cos\theta_2-\cos\theta_1)(\cos\theta_2+\cos\theta_1)$，
非零碰撞时间对应
\[
 k=\frac{2v_0\sin(\theta_1-\theta_2)}
 {g(\cos\theta_2-\cos\theta_1)(\cos\theta_2+\cos\theta_1)}.
\]
从而两次发射间隔
\[
 \boxed{\Delta=t_1-t_2=
 \frac{2v_0\sin(\theta_1-\theta_2)}
 {g(\cos\theta_1+\cos\theta_2)}\simeq10.9\,\mathrm s.}
\]
\excheck 将以上 $t_1,t_2$ 代回水平和竖直位置式，两式均相等；
且 $t_1>t_2>0$，故确为后发射炮弹追上前弹的碰撞。


\exproblem{34}{6}
一气球自地面起以恒定竖直速度 $v_0>0$ 上升；
受风影响，其水平速度为 $v_x=ay$，其中 $y$ 是升空高度，$a$ 是常数。
求：(a) 水平漂移 $x(y)$；(b) 总加速度、切向加速度和法向加速度的大小随 $y$ 的变化。
\exsolution
竖直速度 $\dot y=v_0$ 恒定，故
\[
 \frac{\mathrm dx}{\mathrm dy}=\frac{\dot x}{\dot y}=\frac{ay}{v_0}
 \quad\Longrightarrow\quad
 \boxed{x(y)=\frac{ay^2}{2v_0}},
\]
其中取 $x(0)=0$。瞬时速度、加速度为
$\bm v=(ay,v_0)$、$\bm a=(av_0,0)$，
从而
\[
 \boxed{|\bm a|=|a|v_0,\quad
 a_\tau=\frac{a^2yv_0}{\sqrt{a^2y^2+v_0^2}},\quad
 a_n=\frac{|a|v_0^2}{\sqrt{a^2y^2+v_0^2}}.}
\]
这里 $a_\tau$ 为上升阶段的非负切向加速度。
\excheck $a_\tau^2+a_n^2=(av_0)^2$；
$y=0$ 时速度竖直、加速度水平，故 $a_\tau=0$。

\exproblem{35}{6}
质点速度为 $\bm v=a\bm i+bx\bm j$，
$\bm i,\bm j$ 为坐标轴单位矢量，$a,b$ 为常数，
初始位置 $x=y=0$。求：(a) 轨迹 $y(x)$；
(b) 曲率半径 $R(x)$。
\exsolution
当 $a\ne0$，
\[
 \dot x=a,\quad\dot y=bx,\quad
 \frac{\mathrm dy}{\mathrm dx}=\frac{bx}{a}
 \quad\Longrightarrow\quad
 \boxed{y(x)=\frac b{2a}x^2}.
\]
对显函数曲线 $y(x)$，曲率半径为
$R=(1+y'^2)^{3/2}/|y''|$，
所以在 $b\ne0$ 时
\[
 \boxed{R(x)=\left|\frac ab\right|
 \left(1+\frac{b^2x^2}{a^2}\right)^{3/2}.}
\]
\excheck 直接计算加速度 $\bm a=(0,ab)$，
再用 $R=|\bm v|^3/|\bm v\times\bm a|$ 得到同式。
若 $b=0$ 则轨迹为直线、$R=\infty$，公式需按极限解释。

\exproblem{36}{6}
质点沿一条已给定的轨迹单向运动。
设一恒定矢量 $\bm a_0=a\,\bm i$（$a>0$），
实际切向加速度是 $\bm a_0$ 在单位切向 $\bm\tau$
上的投影 $a_\tau=\bm a_0\cdot\bm\tau$（原书图1.4）。
若经过 $x=0$ 时速度可视为零，求速率 $v(x)$。
\exsolution
定义弧长 $s$，则 $v=\mathrm ds/\mathrm dt$、
$a_\tau=v\,\mathrm dv/\mathrm ds$，
且 $\bm\tau\cdot\bm i=\mathrm dx/\mathrm ds$。因此
\[
 v\frac{\mathrm dv}{\mathrm ds}
 =a\frac{\mathrm dx}{\mathrm ds}
 \quad\Longrightarrow\quad
 \frac12\mathrm d(v^2)=a\,\mathrm dx .
\]
将 $v(0)=0$ 作为初始条件积分：
\[
 \boxed{v(x)=\sqrt{2ax}}\quad (x\ge0).
\]
\excheck 结果与轨迹具体形状无关，但只适用于质点确能到达的
$x\ge0$ 区域；原书的简答因根号字体可能易误认，需保留根号。

\exproblem{37}{6--7}
质点沿圆周运动，速率 $v=at$，
$a=0.50\,\mathrm{m/s^2}$。求从静止开始走过圆周长度的
$n=0.10$ 倍时的合加速度大小。
\exsolution
切向加速度为 $a_\tau=\dot v=a$；
圆周弧长 $s(t)=\frac12at^2$。
走过 $n$ 圈对应 $s=2\pi Rn$，
故 $t^2=4\pi Rn/a$。
法向加速度
\[
 a_n=\frac{v^2}{R}=\frac{a^2t^2}{R}=4\pi na .
\]
因此
\[
 \boxed{a_{\rm total}=a\sqrt{1+16\pi^2n^2}
 \simeq 0.80\,\mathrm{m/s^2}.}
\]
\excheck $R$ 最终消去；$n\to0$ 时合加速度趋于 $a$。

\exproblem{38}{7}
质点沿半径 $R$ 的圆周减速运动，在任意时刻其切向与法向
加速度大小相同，初速为 $v_0$。求：
(a) $v(t)$ 和 $v(s)$；(b) 合加速度大小作为 $v$ 或已走路程 $s$ 的函数。
\exsolution
由于正沿轨迹减速，$a_\tau=\dot v=-v^2/R$。
分离变量：
\[
 \int_{v_0}^{v(t)}\frac{\mathrm dv}{v^2}=-\frac tR
 \quad\Longrightarrow\quad
 \boxed{v(t)=\frac{v_0}{1+v_0t/R}}.
\]
又 $\dot v=v\,\mathrm dv/\mathrm ds$，所以
$\mathrm dv/\mathrm ds=-v/R$，得
\[
 \boxed{v(s)=v_0e^{-s/R}}.
\]
合加速度为互相垂直的两分量的平方和开方：
\[
 \boxed{a_{\rm total}(v)=\frac{\sqrt2\,v^2}{R},\qquad
 a_{\rm total}(s)=\frac{\sqrt2\,v_0^2}{R}e^{-2s/R}.}
\]
\excheck 切向减速始终为负，$t\to\infty$ 时速率趋零，
但不会在有限时刻突然停止；$s\to\infty$。

\exproblem{39}{7}
质点沿半径为 $R$ 的圆弧运动，速率随已走弧长 $s$ 变化为
$v=a\sqrt s$，$a$ 为正常数。求合加速度与速度方向的夹角 $\theta(s)$。
\exsolution
切向加速度 $a_\tau=v\,\mathrm dv/\mathrm ds
 =(a\sqrt s)(a/(2\sqrt s))=a^2/2$；
法向加速度 $a_n=v^2/R=a^2s/R$。二者互相垂直，
故在 $s>0$ 时
\[
 \boxed{\tan\theta=\frac{a_n}{a_\tau}=\frac{2s}{R},
 \qquad \theta(s)=\arctan\frac{2s}{R}.}
\]
\excheck 当 $s\to0^+$，角度趋零；
当 $s\gg R$ 时，法向分量占主导，角度趋 $\pi/2$。

\exproblem{40}{7}
质点沿半径为 $R$ 的圆弧运动，其相对起点的有符号弧位移为
$\ell(t)=a\sin\omega t$。给定 $R=1.00\,\mathrm m$、
$a=0.80\,\mathrm m$、$\omega=2.00\,\mathrm{rad/s}$。
求：(a) $\ell=0$ 及 $\ell=\pm a$ 时合加速度大小；
(b) 最小合加速度及相应弧位移。
\exsolution
沿圆弧的速度与切向加速度为
$\dot\ell=a\omega\cos\omega t$、
$a_\tau=\ddot\ell=-\omega^2\ell$。
法向加速度大小
$a_n=(\dot\ell)^2/R=\omega^2(a^2-\ell^2)/R$。
因此
\[
 A^2(\ell)=\omega^4\left[
 \ell^2+\frac{(a^2-\ell^2)^2}{R^2}\right].
\]
当 $\ell=0$ 时 $A=\omega^2a^2/R=2.56\,\mathrm{m/s^2}$；
$\ell=\pm a$ 时 $A=\omega^2a=3.20\,\mathrm{m/s^2}$。
置 $z=\ell^2\in[0,a^2]$ 并对 $A^2$ 求导：
\[
 \frac{\mathrm d}{\mathrm dz}
 \left[z+\frac{(a^2-z)^2}{R^2}\right]
 =1-\frac{2(a^2-z)}{R^2}.
\]
本题 $a^2-R^2/2=0.14\,\mathrm{m^2}>0$，
内点极小值发生于
$\ell_m=\pm\sqrt{a^2-R^2/2}$。故
\[
 \boxed{\ell_m\simeq\pm0.374\,\mathrm m,\qquad
 A_{\min}=\omega^2\sqrt{a^2-\frac{R^2}{4}}
 \simeq2.50\,\mathrm{m/s^2}.}
\]
\excheck 计算表明 $2.50<2.56<3.20$。
一般参数下若 $a^2\le R^2/2$，极小点在 $\ell=0$，
不能不核查区间就套内点公式。
